\subsection{Actual affine-norm conjugacy at the quadratic and IV caps}

\begin{statement}{P3F.49}{Proved}
At the quadratic $I_1^*$ root and the good or reduced IV roots of
P3F.36 and P3F.38, every extending affine source norm solution
with the prescribed canonical deck and the permitted transitive
central component action is conjugate to the constructed solution
by an extending translation commuting with the suspension.
For the actual multidegree-zero root line, this conjugacy lifts
to an invertible line map. After that lift, the remaining norm
solution differs by a retained finite delta character.
The actual outer collar coefficients remain in the comparison.
\end{statement}
\begin{proof}
\par\medskip\noindent\textbf{The translation equation.}
Let $D$ be the canonical order-$m$ root deck. Write the constructed
and the given source maps as
\[
 A_0=t_{R_0}D,\qquad A=t_R D,\qquad
 A^m=A_0^m=t_{-nQ}.
\]
Here $R,R_0$ are holomorphic sections of the root Neron smooth
group. Its group law is written additively, except in the explicit
multiplicative central calculation below. The action of $D$ on
sections is the one defined by $Dt_TD^{-1}=t_{DT}$.
Then $S=R-R_0$ satisfies
\[
 \operatorname{Nm}_D(S)=\sum_{j=0}^{m-1}D^jS=0.               \tag{F49.1}
\]
Conjugation by the actual translation $t_T$ changes $R_0$ to
$R_0+(1-D)T$. We solve $(1-D)T=S$ on the whole root disc.

For a good IV root, the central group is the smooth elliptic
curve $E_0$, and $D$ acts with order three. The actual isogeny
$1-D:E_0\to E_0$ has degree three and is onto. Choose a central
solution $T_0$ and extend it to a local holomorphic section by
the smooth submersion coordinate of the root family.

For a reduced IV root, $D$ fixes the component group
$\mathbb Z/3$ and $m\equiv1\pmod3$. Equation (F49.1) therefore
puts $S_0$ in the additive identity component. If $\lambda$ is
the scalar action of $D$ on that component and $\lambda\ne1$,
take $T_0=S_0/(1-\lambda)$. If $\lambda=1$, the central norm
equation is $mS_0=0$, so $S_0=0$ and we take $T_0=0$.
Again the smooth group submersion extends this central section.

For the quadratic starred root, $m=2$. Both allowed source
maps exchange the two $I_2$ components, so both translation
components are odd. Their difference lies in the multiplicative
identity component $\mathbb C^*$. On it $D$ is inversion, so
the actual homomorphism $1-D$ is
\[
                     T\longmapsto T^2.
\]
Choose $T_0$ with $T_0^2=S_0$ and extend it in the identity
component. This calculation retains the transitive component
class required by the cap construction.

Subtracting $(1-D)T_0$ leaves a norm-zero section whose central
value is the identity. Shrink the root disc so that it lies in
the analytic formal-group neighbourhood of the identity. Its
fibrewise logarithm is a section $s$ of the Lie line. The action
of $D$ on this line is linear over its finite base rotation.
The norm of $s$ is lattice valued and vanishes at the centre;
after the same shrinking it is zero. The explicit finite sum
\[
 t=-\frac1m\sum_{j=0}^{m-1}jD^js,
 \qquad
 (1-D)t=s-\frac1m\sum_{j=0}^{m-1}D^js=s                    \tag{F49.2}
\]
solves the residual equation. Exponentiate and add this
translation to the first one. We have obtained the actual maps
\[
             t_TA_0t_T^{-1}=A,
 \qquad t_Tt_{-Q}=t_{-Q}t_T.                               \tag{F49.3}
\]
The translation extends on the minimal root surface. On a good
root it is the smooth elliptic group action. On the $I_2$ root
it is multiplication by the corresponding Neron unit point.
On the reduced IV root it is the extending section translation
in the resolved charts used in P3F.38. Thus (F49.3) is an
identity of the actual extending source maps.

The quadratic case also has a literal Tate-unit formula. With
the prescribed boundary marking, write
\[
 A(Z,z)=(R(z)/Z,-z),\quad R(z)=z w(z),\qquad
 R_0(z)=z w_0(z).
\]
Equality of the two actual norms makes $S=w/w_0$ an even
nowhere-zero unit. Choose its even analytic square root $U$.
Then $C_U(Z,z)=(U(z)Z,z)$ gives
\[
 C_UA_0C_U^{-1}=A,
 \qquad U(z)U(-z)R_0(z)=R(z).                              \tag{F49.3a}
\]
This map can move the four canonical inversion centres.
We transport the Cartier comparison by the literal blowup
isomorphism
$\operatorname{Bl}_{Z_0}S'\to
\operatorname{Bl}_{C_U(Z_0)}S'$ induced by $C_U$; the new
quotient map is the old quotient map composed with its inverse.
Thus the exceptional divisor comparison is transported, without
requiring $C_U$ to fix its original four centres.

\par\medskip\noindent\textbf{The line lift and its finite phase.}
Let $N$ be the actual root line, including the Cartier corrections
proved in P3F.36 and P3F.38. It has multidegree zero. Translation
by $T$ preserves its generic fibrewise Picard class. Apply the
proper-direct-image and vertical-degree calculation of P3K.15
to the ratio of $N$ and its translate. Evaluation gives a base
line, and the remaining divisor is an integral whole fibre.
Its local base equation supplies an actual invertible lift
$L_T$; we retain that base factor and its outer phase.

These degree-zero translation lifts commute with the specified
suspension $B$. On every smooth fibre, write the punctured line
as $(\mathbb C\times\mathbb C^*)/\Lambda$ with constant
multiplicative period characters. A translation lift acts by
a constant translation in the first factor and a constant scalar
in the second. The two such maps commute. Every other lift
differs by a base scalar, so this proves $[L_T,B]=1$ on the
generic fibres and, by extension, on the root neighbourhood.

The conjugated canonical lift $L_TA_0L_T^{-1}$ and the given
lift of $A$ now cover the same source map. Their ratio is a
holomorphic base unit. More precisely, if
$C=L_TA_0L_T^{-1}$, define $f$ by $A=m_fC$, where $m_f$
multiplies the vector in its \emph{target} fibre by $f$.
Thus conjugation by a base scalar $b$ has target ratio
$b(z)/b(\zeta_m^{-1}z)$, in agreement with the action of $D$
used in (F49.2). Their common $m$th power $B^n$
gives the exact norm equation
\[
               \prod_{j=0}^{m-1}f(\zeta_m^jz)=1.             \tag{F49.4}
\]
Put $\chi=f(0)$, so $\chi^m=1$. The unit $f/\chi$ has an
analytic logarithm $g$ with $g(0)=0$, and its norm is zero.
Formula (F49.2), applied to this scalar logarithm with the same
conjugation action, removes $\exp(g)$ by a nowhere-zero base
scalar $b$, explicitly
$A=\chi\,m_b C m_b^{-1}$. For $m=2$, $g$ is odd and
$b=\exp(g/2)$ has the required target ratio
$b(z)/b(-z)=\exp(g)$. The remaining constant $\chi$ is retained as the finite
delta character of the cap.

Consequently the line comparison carries the exact norm maps,
the finite character, and all outer source and line markings.
The homology calculations of P3F.36 and P3F.38 already retain
their actual coefficients $c,\kappa$; no choice of a numerical
outer coefficient is needed here. The displayed conjugacy is
local at each cap and makes no change between distinct source
sites. The exact averaging and central squaring identities are
recorded in \path{code/part3/os50_affine_conjugacy.py} and its
executed JSON output.
\end{proof}
